\begin{abstract}
Enterprise data platforms track tens of thousands of tables. Their technical metadata
(schemas, lineage) is captured automatically, but their semantic metadata (what a table
means, what one row represents, what each column measures) is usually missing, which slows
data discovery, governance, and analytics. We present OpenMetaDataGenerator (OMDG), an
open-source system that generates this semantic layer by grounding a large language model
(LLM) in three signals: the table schema, external context (transformation code and
documentation), and data lineage. OMDG contributes (i) a topological wave scheduler that
generates descriptions in lineage order, so downstream tables inherit the already-generated
descriptions of their upstreams; (ii) canonicalize-first generation, which describes each
recurring column concept once and reuses it across the catalog; (iii) a controlled agentic
loop that measures coverage and per-item accuracy against explicit targets and selects the
next gap-filling strategy (inherit, sibling, or rework), with an LLM judge deciding whether
a reworked description replaces the previous one; and (iv) an auditable substrate consisting
of a local DuckDB knowledge store, a persistent incremental HNSW retrieval index, and a
confidence tag on every description. The system supports DataHub, Databricks, and Snowflake
as metadata sources and six LLM providers. Because production catalogs lack ground-truth
semantics, we evaluate on a labelled synthetic benchmark, two public schemas (TPC-H and
Sakila) whose foreign keys we read as lineage, and a cryptic-name variant of Sakila with
obfuscated identifiers. With a strong production model (Gemini~3.1~Pro) we observe two
regimes. When column names are self-describing, the model is already accurate from names
alone and context adds at most a few points. When names carry no meaning, grounding
determines the outcome: accuracy rises from 0.27 to 0.71 ($+43$ points, stable across three
runs) on the cryptic variant, whose gold labels and lineage are identical to the readable
one. The benefit of grounding is therefore governed by name informativeness, and it is
largest for the cryptic, abbreviated columns that are common in enterprise catalogs. We
specify a human-evaluation protocol for production deployments and release the system,
benchmarks, and evaluation harness under Apache-2.0.
\end{abstract}

\section{Introduction}
The value of a data platform depends on whether its users can find and trust the right
data. Technical metadata (column names, types, table lineage) is now captured automatically
by catalog systems~\citep{halevy2016goods,hellerstein2017ground}. Semantic metadata, that
is, a one-sentence statement of what a table is, its grain (what a single row represents),
and the meaning of each column, remains largely manual and is missing for the long tail of
tables. In practice, analysts re-derive tables that already exist, governance teams cannot
classify sensitive columns, and downstream contracts drift from reality.

Large language models~\citep{vaswani2017attention,brown2020language} make automatic
description generation feasible, but prompting a model with only a table's column list
produces generic and often hallucinated text. What such prompts lack is grounding: tables
are produced by transformation code, are partially described in scattered documentation,
and are connected to other tables by lineage that carries meaning. We argue that reliable
descriptions require combining these signals, and that lineage should also determine the
order in which tables are documented, not only the content of individual prompts.

\paragraph{Contributions.}
\begin{enumerate}
  \item \textbf{Topological wave scheduling.} Generation proceeds in topological (lineage)
  order, so when a downstream table is described, the already-generated descriptions of its
  upstream tables are available as context. Grounding therefore propagates through the
  lineage DAG (Section~\ref{sec:method}).
  \item \textbf{Retrieval-grounded, provenance-tagged prompting.} Each prompt combines
  schema signals, retrieved transformation code and documentation, and fine-grained column
  lineage; every output records which signals were used
  (Sections~\ref{sec:method}--\ref{sec:system}).
  \item \textbf{Canonicalize-first generation.} The same logical column recurs across many
  tables in an enterprise schema. We normalize column names, generate one description per
  frequent concept, and reuse it everywhere, which reduces cost and keeps wording consistent.
  \item \textbf{A controlled agentic loop.} Instead of a fixed pipeline, a strategy
  controller scores the whole catalog against coverage and accuracy targets in a single
  batched similarity evaluation, then selects the next action through a short, logged chain
  of thought over a constrained action set: inherit (copy a description from a same-named or
  upstream column), sibling (infer a table's empty columns from its described ones), rework
  (regenerate the lowest-similarity items with a rising sampling temperature), or stop. An
  LLM judge decides whether a reworked description replaces the previous one, and every
  decision is logged with its rationale, temperature, and measured state. The loop stops
  when the targets are met or no strategy yields further improvement.
  \item \textbf{An auditable, incremental substrate.} A local DuckDB store persists
  metadata, table- and column-level lineage, and each description's grounding context and
  similarity score. A model-stamped HNSW index over code, documentation, and view-SQL
  chunks makes retrieval persistent and incremental: only changed inputs are re-embedded.
  Every description carries a confidence tag so review effort can be directed at the weakly
  grounded tail.
  \item \textbf{A labelled benchmark suite and open system.} We release a synthetic
  benchmark, two public-schema benchmarks (TPC-H, Sakila), a cryptic-name benchmark, an
  evaluation harness (coverage, accuracy, faithfulness), and a source- and model-agnostic
  reference implementation.
\end{enumerate}

\section{Related Work}
OMDG builds on three lines of work: data cataloging, LLMs for data management, and
retrieval-augmented generation.

\paragraph{Data cataloging and context services.} Systems such as Goods inventory an
organization's datasets and their provenance~\citep{halevy2016goods}, and Ground models the
surrounding data context~\citep{hellerstein2017ground}. They automate the capture of
technical metadata but leave semantic description to humans. OMDG targets that layer and
writes its output back into such catalogs.

\paragraph{LLMs for data management.} Foundation models have been applied to data wrangling
and entity matching~\citep{narayan2022foundation} and to synthesizing query code from
natural language~\citep{trummer2022codexdb}. These works apply LLMs to transform or query
data. OMDG applies them to document data, and conditions generation on the catalog's
lineage structure rather than treating each table in isolation.

\paragraph{Semantic column annotation and table understanding.} A related line of work
classifies columns into fixed semantic-type vocabularies: Sherlock and Sato use feature-
and context-based models~\citep{hulsebos2019sherlock,zhang2020sato}, and DODUO uses
pre-trained language models~\citep{suhara2022doduo}. TURL learns table representations for
downstream tasks~\citep{deng2020turl}. Classification into a closed type set supports
matching and governance but cannot express a table's grain or business meaning. OMDG
produces free-text descriptions grounded in code, documentation, and lineage, which is a
more expressive and harder-to-evaluate target.

\paragraph{Retrieval-augmented generation and its evaluation.} Grounding generation in
retrieved evidence is standard practice~\citep{lewis2020rag}, usually over a single text
corpus with dense sentence-embedding retrieval~\citep{reimers2019sentencebert}. Recent work
evaluates such systems for faithfulness without gold references~\citep{es2024ragas}, which
motivates our context-versus-output similarity proxy. Our retrieval covers several sources
(transformation code, documentation, view SQL) and is combined with a structural signal:
the lineage DAG determines generation order, so retrieved and generated upstream context
propagates to downstream tables. Generic retrieval-augmented generation does not use graph
structure in this way.

\section{Problem Formulation}
Let a catalog be a set of tables $T=\{t_1,\dots,t_n\}$. Each table $t$ has columns
$C(t)$, an optional (usually empty) existing description, an optional view/SQL definition,
and a lineage relation: coarse upstreams $U(t)\subseteq T$ and, per column $c$,
fine-grained upstreams $U(c)$. External corpora provide code $\mathcal{K}$ and
documentation $\mathcal{D}$. The goal is a function $g$ mapping each table and column to a
description that is correct (faithful to the table's true semantics) and grounded
(supported by available evidence). We evaluate against gold descriptions $y^\star$ using
semantic similarity, and measure hallucination via the similarity of an output to the
context it was conditioned on.

\section{Method}\label{sec:method}
\paragraph{Topological wave scheduling.} We build a dependency graph over the in-scope
tables from $U(\cdot)$ and partition it into waves with Kahn's algorithm
(Algorithm~\ref{alg:main}); cycles, which are rare (for example mutually recursive views),
are broken by emitting the remaining nodes as a final wave. Tables in wave $k$ are
generated only after all earlier waves, so the prompt for a downstream table includes the
freshly generated descriptions of its upstreams.

\paragraph{Retrieval-grounded prompting.} For each table $t$ the prompt combines five
signals: the schema layer (a medallion prior inferred from the schema name), the typed
column list, the view SQL when present, the top-$k$ chunks retrieved from the code and
documentation corpora $\mathcal{K},\mathcal{D}$ (exact-name matches plus dense retrieval
over the persistent index of Section~\ref{sec:system}), the inherited descriptions of
$t$'s upstreams, and fine-grained column-lineage hints. The model returns JSON with a table
description and per-column descriptions, and is instructed to describe only columns it can
ground in the supplied evidence. Columns with insufficient evidence are left empty for the
control loop instead of being guessed.

\paragraph{Canonicalize-first.} Before wave generation, we canonicalize column names
(normalizing case and separators so that \texttt{Order\_ID}, \texttt{order\_id}, and
\texttt{orderid} collapse to one key) and, for every concept that recurs in at least $m$
tables with no table-specific lineage, generate a single description and copy it onto all
occurrences. Wave generation then only has to describe the remaining table-specific
columns. This reduces cost (one call replaces many) and keeps wording identical for
identical concepts.

\paragraph{Controlled agentic loop.} After the wave pass and a table-level coverage fill,
control passes to a strategy controller (Algorithm~\ref{alg:main}). Each iteration has four
steps. (1)~Similarity evaluation: every described object is scored against the exact
context it was grounded in, $s(x)=\cos(\phi(\hat y_x),\phi(\mathrm{ctx}_x))$, in one
batched encoder pass; the mean is the current accuracy and the low-scoring items are the
rework candidates. (2)~Decision: given the measured coverage, accuracy, and per-strategy
candidate counts relative to the targets $\tau_{\text{cov}},\tau_{\text{acc}}$, the LLM
reasons in one short sentence and commits to a single action from a constrained set, a
bounded and logged form of chain-of-thought prompting~\citep{wei2022chain}, with a
deterministic rule (fill coverage before accuracy) as fallback. (3)~Action: inherit copies
a description from a same-named or fine-grained-upstream column at no model cost; sibling
infers a table's still-empty columns from its described ones with one small call per table;
together these fill exactly the undescribed objects toward $\tau_{\text{cov}}$. Rework
regenerates the lowest-similarity items with a vocabulary-reuse instruction, raising the
sampling temperature on each successive pass to escape local minima. Stop ends the loop. A
rework output replaces the current description only if it differs and an LLM judge rules it
an improvement. (4)~Bookkeeping: coverage strategies run at most once; rework recurs while
it yields accepted improvements. The loop terminates when the targets are met, no strategy
yields further change, or an iteration cap is reached. Each iteration's action, rationale,
temperature, and measured state are appended to a decision trace, so the process is bounded
and auditable.

\paragraph{Confidence tagging.} Every emitted description is prefixed with a confidence
tag: \texttt{[AIG\,|\,High]} when the object was grounded in code, documentation, or
lineage; \texttt{[AIG\,|\,Low]} when only the schema and column names were available; and
\texttt{[Reviewed]} is preserved when a human has edited a description. The tag travels
with the description into the catalog, so consumers see a confidence signal and human
review can be prioritized on the low-confidence tail. Tags are applied after scoring, so
they do not affect the accuracy metric.

\begin{algorithm}[t]
\caption{OMDG: canonicalize-first, lineage-aware, controlled agentic generation}\label{alg:main}
\begin{algorithmic}[1]
\Require tables $T$; LLM $M$; judge $J$; embedder $\phi$; targets $\tau_{\text{cov}},\tau_{\text{acc}}$; threshold $\theta$
\State \textsc{CanonicalizeFirst}$(T,M)$ \Comment{one desc per frequent concept, seeded everywhere}
\State $W \gets \textsc{TopoWaves}(T)$; \; $D \gets \{\}$ \Comment{Kahn over in-scope lineage}
\For{wave $w \in W$} \Comment{upstream-first so grounding propagates}
  \State $\{\hat y_t\} \gets M.\textsc{Batch}(\{\textsc{Prompt}(t,D):t\in w\})$; update $D$
\EndFor
\State \textsc{CoverageFill}$(T,D)$; \; $E \gets \emptyset$ \Comment{$E$: exhausted strategies}
\While{iterations left}
  \State $g \gets \textsc{MeasureGaps}(T)$ \Comment{coverage, accuracy, per-strategy candidates}
  \State $a \gets \textsc{Decide}(g, E, \tau_{\text{cov}}, \tau_{\text{acc}})$ \Comment{LLM + deterministic fallback}
  \If{$a=\textbf{stop}$} \textbf{break} \EndIf
  \State $k \gets \textsc{Apply}(a, T, M, J)$ \Comment{inherit/sibling/rework; rework judged by $J$}
  \If{$a\in\{\text{inherit},\text{sibling}\}$ or $k=0$} $E \gets E \cup \{a\}$ \EndIf
\EndWhile
\State \Return descriptions with per-object provenance and the decision trace
\end{algorithmic}
\end{algorithm}

\section{System}\label{sec:system}
OMDG is organized around pluggable interfaces: a metadata source (DataHub via GraphQL;
Databricks and Snowflake via \texttt{information\_schema}), context providers (a code path
and a document path), and an LLM backend (OpenAI, Anthropic, AWS Bedrock, Google Vertex,
Azure OpenAI, and the Gemini API behind one interface). Between the source and the
generator sits a local DuckDB store~\citep{raasveldt2019duckdb} with relations
\texttt{kb\_tables}, \texttt{kb\_columns}, and \texttt{kb\_embed\_index}. It persists
pulled metadata, table- and column-level lineage, the exact \texttt{prompt\_context} and
per-item similarity of each description, and the embedding vectors of the retrieval corpus.
Retrieval uses a persistent, incremental HNSW index~\citep{malkov2020hnsw}: context chunks
(code, documentation, view SQL) are embedded once with a local sentence encoder, stamped
with the model name, and re-embedded only when their source changes; a model change wipes
and rebuilds the index. Persisting this state makes runs incremental (only changed objects
are re-documented) and auditable end to end. The data flow is
\texttt{source}$\to$\texttt{DuckDB (store + HNSW retrieval)}$\to$\texttt{LLM (agentic
generation)}$\to$\texttt{write-back}. The final export is a CSV with one row per described
object: \texttt{(object\_type, object\_name, parent\_table, context, output,
grounding\_notes, similarity)}. Pairing each output with the exact context it saw makes the
export directly usable by an external judge or faithfulness pipeline~\citep{es2024ragas}.

\section{Experiments}
\paragraph{Synthetic benchmark.} Real catalogs cannot supply ground truth, so we generate a
synthetic catalog from a library of domain concepts (each with a known grain and typed,
described columns) arranged in a medallion DAG
(\texttt{raw}$\to$\texttt{intermediate}$\to$\texttt{mart}). Every table receives gold table
and column descriptions; code and documentation context are granted per table with
probability $p$, which lets us ablate context availability. Lineage inheritance is ablated
by severing $U(\cdot)$ before generation.
\paragraph{Metrics.} Coverage (fraction of objects described), accuracy (mean
sentence-embedding cosine of output versus gold), and faithfulness (mean cosine of output
versus the grounding context, a hallucination proxy). The harness also computes grain
recovery, but we do not tabulate it: exact-phrase matching mis-scores the varied phrasings
an LLM uses for the same grain (``each row represents \ldots'' versus ``one row per
\ldots''), so grain is assessed in the human protocol instead.
\paragraph{Public-schema benchmarks (TPC-H, Sakila).} To complement the synthetic setting,
we evaluate on two public schemas whose names, types, and foreign-key relationships we did
not author. TPC-H (8 tables) forms a dependency DAG rooted at \texttt{region} and
\texttt{nation} and flowing through \texttt{supplier} and \texttt{customer} down to
\texttt{orders} and \texttt{lineitem}. Sakila (15 tables) is larger and more connected: it
contains a \texttt{store}--\texttt{staff} cycle, which exercises cycle handling, and a
\texttt{last\_update} column present in every table, which exercises canonicalize-first
(the concept is described once and copied to all 15 occurrences). Foreign keys are read as
lineage edges; gold descriptions are paraphrased from each schema's public documentation,
and documentation context is withheld or limited so that grounding must come from schema
and lineage rather than from a leaked target.

\paragraph{Cryptic-name benchmark.} The setting that motivates OMDG is a catalog whose
identifiers carry little meaning. To measure it with public data, we take Sakila, keep its
gold descriptions, foreign keys, and DAG, and deterministically obfuscate every identifier
(\texttt{customer} $\to$ \texttt{t05}, \texttt{customer\_id} $\to$ \texttt{c05\_00}), which
decouples name from meaning. The grounding context is realistic and does not leak the
answer: it consists of the rename ETL that built each warehouse table from a readable
legacy source (\texttt{SELECT customer\_id AS c05\_00, \ldots\ FROM legacy.customer}) plus
a one-line catalog note. The gold description text never appears in any context, so the
model must still know what the mapped source column means. Lineage, including fine-grained
FK$\to$PK column lineage, is preserved under the renaming.

\paragraph{Results.} Tables~\ref{tab:main}--\ref{tab:cryptic} report the four
configurations generated with Gemini (\texttt{gemini-3.1-pro-preview}, the model our
reference deployment uses) and scored by sentence-embedding cosine. The experiments show
two regimes, separated by how informative the schema names are.
With descriptive names (TPC-H and the synthetic set), a strong production model is already
accurate from column names alone and context adds nothing measurable (within $\pm1$ point).
On Sakila, the largest of the readable schemas, context adds $+5.4$ points, and the
configurations order as full $>$ no-lineage $>$ no-context $>$ schema-only (single run
each).
With cryptic names (Table~\ref{tab:cryptic}), grounding determines the outcome. Mean
accuracy over three independent runs rises from $0.273$ without context to $0.707$ with it,
a gain of $43.4$ points ($2.6\times$), with run-to-run standard deviation at most $0.006$
per configuration. Without any grounding the model also declines objects it cannot
describe (schema-only coverage is $0.91$ in all three runs); context restores full
coverage. Because the schema, gold labels, and lineage are identical across the two Sakila
variants and only the names differ, the comparison is controlled and the gap is
attributable to name informativeness. Lineage inheritance shows no consistent accuracy
effect in either regime ($\pm1$--$2$ points); its value is coverage and cross-table
consistency, which a capable model already saturates on these schemas.
A deterministic backend that requires no API keys reproduces the pipeline and its ablation
mechanics for anyone without credentials; because it does not semantically exploit
inherited descriptions, it is lineage-neutral by construction, and the numbers above are
the ones to cite for lineage's small effect.

\begin{table}[t]
\centering
\small
% auto-generated by benchmark/run.py
\begin{tabular}{lccc}
\toprule
Configuration & Coverage & Accuracy & Faithfulness \\
\midrule
Schema-only (baseline) & 1.00 & 0.672 & 0.416 \\
-- code/doc context & 1.00 & 0.726 & 0.380 \\
-- lineage inheritance & 1.00 & 0.676 & 0.431 \\
Full (context + lineage) & 1.00 & 0.723 & 0.390 \\
\bottomrule
\end{tabular}

\caption{Synthetic benchmark (5 concepts, 15 tables, 80 objects); Gemini
\texttt{gemini-3.1-pro-preview}, embedding-cosine, single run. The model recovers the
templated targets from column names alone, so context and lineage have negligible effect
here. Coverage is 1.0 throughout.}
\label{tab:main}
\end{table}

\begin{table}[t]
\centering
\small
% auto-generated by benchmark/run.py
\begin{tabular}{lccc}
\toprule
Configuration & Coverage & Accuracy & Faithfulness \\
\midrule
Schema-only (baseline) & 1.00 & 0.805 & 0.267 \\
-- code/doc context & 1.00 & 0.798 & 0.299 \\
-- lineage inheritance & 1.00 & 0.826 & 0.324 \\
Full (context + lineage) & 1.00 & 0.753 & 0.335 \\
\bottomrule
\end{tabular}

\caption{Public TPC-H schema (8 tables, 66 objects); foreign keys read as lineage. The
model is already accurate from names alone: context and lineage change accuracy by less
than one point.}
\label{tab:tpch}
\end{table}

\begin{table}[t]
\centering
\small
% auto-generated by benchmark/run.py
\begin{tabular}{lcccc}
\toprule
Configuration & Coverage & Accuracy & Faithfulness & Grain \\
\midrule
Schema-only (baseline) & 1.00 & 0.245 & 0.133 & 0.00 \\
-- code/doc context & 1.00 & 0.279 & 0.139 & 0.00 \\
-- lineage inheritance & 1.00 & 0.309 & 0.157 & 1.00 \\
Full (context + lineage) & 1.00 & 0.344 & 0.162 & 1.00 \\
\bottomrule
\end{tabular}

\caption{Public Sakila schema (15 tables, 96 objects); a deeper foreign-key DAG with a
\texttt{store}--\texttt{staff} cycle and a catalog-wide \texttt{last\_update} column. The
only readable schema where context helps ($+5.4$ accuracy points); the configurations order
as schema-only $<$ no-context $<$ no-lineage $<$ full.}
\label{tab:sakila}
\end{table}

\begin{table}[t]
\centering
\small
% auto-generated: mean +/- sd over 3 independent runs (merge_cryptic.py)
\begin{tabular}{lccc}
\toprule
Configuration & Coverage & Accuracy & Faithfulness \\
\midrule
Schema-only (baseline) & 0.91$\pm$0.00 & 0.291$\pm$0.003 & 0.308$\pm$0.003 \\
-- code/doc context & 1.00$\pm$0.00 & 0.256$\pm$0.006 & 0.412$\pm$0.012 \\
-- lineage inheritance & 1.00$\pm$0.00 & 0.711$\pm$0.004 & 0.315$\pm$0.000 \\
Full (context + lineage) & 1.00$\pm$0.00 & 0.704$\pm$0.001 & 0.358$\pm$0.003 \\
\bottomrule
\end{tabular}

\caption{Cryptic-name benchmark: Sakila with obfuscated identifiers (15 tables, 96
objects); mean $\pm$ sd over three independent runs. When names carry no meaning, grounding
determines the outcome: accuracy rises from 0.273 to 0.707 ($+43.4$ points, $2.6\times$),
and without grounding the model declines objects it cannot describe (coverage 0.91).}
\label{tab:cryptic}
\end{table}

\subsection{Human-evaluation protocol}
Embedding similarity to a single reference under-credits correct paraphrases, so we specify
a human study to accompany production deployments (we release the protocol, not proprietary
results). Annotators with domain familiarity rate each generated description on three
1--5 Likert axes: correctness (is every claim true?), completeness (are the grain and key
semantics stated?), and groundedness (is every claim supported by the shown context?). The
study is a blind A/B against the schema-only baseline, with a subset double-annotated to
report inter-annotator agreement (Cohen's $\kappa$). Each item is presented with its
\texttt{grounding\_notes} provenance so raters can verify support, and hallucinations
(unsupported claims) are logged separately as a hard error rate.

\subsection{Reproducibility statement}
The system, the synthetic generator, the public-schema and cryptic-name benchmarks, and the
evaluation harness are released under Apache-2.0 at \REPO. The result tables
(Tables~\ref{tab:main}--\ref{tab:cryptic}) are produced with
\texttt{make benchmark-llm PROVIDER=gemini MODEL=gemini-3.1-pro-preview} and can be
regenerated with any of the six supported providers (absolute numbers vary by model).
Independently, the pipeline and its ablation mechanics run with no API keys and no GPU via
a deterministic reference backend (\texttt{pip install -e .}; \texttt{make benchmark}),
completing in seconds; this path is used for continuous integration and by reviewers
without credentials. Unit tests (\texttt{make test}) and two runnable demos (the controlled
agentic loop and the DuckDB+HNSW retrieval) execute on a fresh install; CI runs the tests
and all benchmarks on Python 3.9--3.12.

\section{Discussion and Limitations}\label{sec:limitations}
\paragraph{The benefit of grounding depends on name informativeness.} On schemas whose
names are descriptive (TPC-H, Sakila, and the templated synthetic set), a strong LLM
already produces accurate descriptions from names alone, so external context adds only a
few points and lineage inheritance adds little to accuracy (it saturates coverage instead).
This is the regime where the task is easy. The cryptic-name benchmark tests the regime OMDG
is built for, columns whose names carry no meaning, by obfuscating identifiers while
holding schema, gold, and lineage fixed; there, grounding is worth $+43$ points. One caveat
is that its grounding context (the rename ETL) contains an explicit name mapping; real
pipelines vary in how directly the mapping can be recovered, so gains in practice will fall
between our two regimes depending on both name quality and context quality.
\paragraph{Runs and metric caveats.} The cryptic-name result is reported as mean $\pm$ sd
over three independent runs (sd at most 0.006); the descriptive-name benchmarks are single
runs, which we consider sufficient for their null findings given the observed run-to-run
stability. Embedding cosine to a single reference under-credits correct paraphrases, and
the rework loop's similarity signal can likewise under-value terse descriptions; pairing
both with a reference-free judge is future work.
\paragraph{Human-in-the-loop.} Generated descriptions should be surfaced with their
grounding provenance and confidence tag so that reviewers stay in the loop for high-stakes
governance decisions rather than trusting the model blindly.

\section*{Broader Impact Statement}
OMDG generates documentation that people may rely on for data discovery and governance. The
primary risk is misplaced trust in incorrect machine-generated descriptions. The system
mitigates this by construction: every description carries a confidence tag and grounding
provenance, ungrounded objects are labelled low-confidence rather than embellished, and the
export pairs each output with the exact evidence it saw, so human review can be targeted at
the least-grounded tail. Deployments should keep a human approval step for descriptions
used in compliance-sensitive settings. The system processes only metadata (schemas, code,
documentation), not record-level data, which limits privacy exposure; operators should
still confirm that the metadata itself is not sensitive before sending it to an external
LLM provider, or use a locally hosted backend.

\section{Conclusion}
Semantic metadata is the missing half of the modern data catalog, and generating it well
requires more than prompting a language model with a column list. OMDG treats lineage as a
scheduling structure rather than only prompt content, and combines it with
canonicalize-first generation, multi-source retrieval grounding, and a bounded, auditable
agentic control loop. Empirically, the benefit of grounding is governed by name
informativeness: on public schemas with descriptive names a strong production LLM is
already a hard baseline and context adds at most a few points, but on the same schema with
obfuscated identifiers grounding raises accuracy from 0.27 to 0.71 ($+43$ points, stable
over three runs). This is direct evidence for the setting that motivates the system: the
cryptic columns of real enterprise catalogs. The design is agnostic to the metadata source
and the LLM provider, persists its state for incremental and auditable runs, and attaches
provenance and confidence to every description so that humans stay in the loop where it
matters. We release the system, the four benchmarks, and the evaluation harness to support
reproducible research on automatic semantic-metadata generation.

% bibliography commands live in each build's main file (styles differ)
